\documentclass[10pt,a4paper]{amsart}
\usepackage{amsmath,amssymb,amsthm,graphicx}
% For commutative diagrams we use Xy-pic
\usepackage[all]{xy}

% We use 2cell for 2-commutative diagrams.
\xyoption{2cell}
% Theorem environments.
%
\theoremstyle{plain}
\newtheorem{theorem}[subsection]{Theorem}
\newtheorem{proposition}[subsection]{Proposition}
\newtheorem{lemma}[subsection]{Lemma}

\theoremstyle{definition}
\newtheorem{definition}[subsection]{Definition}
\newtheorem{example}[subsection]{Example}
\newtheorem{exercise}[subsection]{Exercise}
\newtheorem{situation}[subsection]{Situation}

\theoremstyle{remark}
\newtheorem{remark}[subsection]{Remark}
\newtheorem{remarks}[subsection]{Remarks}

\numberwithin{equation}{subsection}

% Macros
%
\def\lim{\mathop{\mathrm{lim}}\nolimits}
\def\colim{\mathop{\mathrm{colim}}\nolimits}
\def\Spec{\mathop{\mathrm{Spec}}}
\def\Hom{\mathop{\mathrm{Hom}}\nolimits}
\def\Ext{\mathop{\mathrm{Ext}}\nolimits}
\def\SheafHom{\mathop{\mathcal{H}\!\mathit{om}}\nolimits}
\def\SheafExt{\mathop{\mathcal{E}\!\mathit{xt}}\nolimits}
\def\Sch{\mathit{Sch}}
\def\Mor{\mathop{\mathrm{Mor}}\nolimits}
\def\Ob{\mathop{\mathrm{Ob}}\nolimits}
\def\Sh{\mathop{\mathit{Sh}}\nolimits}
\def\NL{\mathop{N\!L}\nolimits}
\def\CH{\mathop{\mathrm{CH}}\nolimits}
\def\proetale{{pro\text{-}\acute{e}tale}}
\def\etale{{\acute{e}tale}}
\def\QCoh{\mathit{QCoh}}
\def\Ker{\mathop{\mathrm{Ker}}}
\def\Im{\mathop{\mathrm{Im}}}
\def\Coker{\mathop{\mathrm{Coker}}}
\def\Coim{\mathop{\mathrm{Coim}}}

% Boxtimes
%
\DeclareMathSymbol{\boxtimes}{\mathbin}{AMSa}{"02}

%
% Macros for moduli stacks/spaces
%
\def\QCohstack{\mathcal{QC}\!\mathit{oh}}
\def\Cohstack{\mathcal{C}\!\mathit{oh}}
\def\Spacesstack{\mathcal{S}\!\mathit{paces}}
\def\Quotfunctor{\mathrm{Quot}}
\def\Hilbfunctor{\mathrm{Hilb}}
\def\Curvesstack{\mathcal{C}\!\mathit{urves}}
\def\Polarizedstack{\mathcal{P}\!\mathit{olarized}}
\def\Complexesstack{\mathcal{C}\!\mathit{omplexes}}
% \Pic is the operator that assigns to X its picard group, usage \Pic(X)
% \Picardstack_{X/B} denotes the Picard stack of X over B
% \Picardfunctor_{X/B} denotes the Picard functor of X over B
\def\Pic{\mathop{\mathrm{Pic}}\nolimits}
\def\Picardstack{\mathcal{P}\!\mathit{ic}}
\def\Picardfunctor{\mathrm{Pic}}
\def\Deformationcategory{\mathcal{D}\!\mathit{ef}}
\usepackage{cleveref}
\crefname{originallemma}{original lemma}{original lemmas}
\Crefname{originallemma}{Original lemma}{Original lemmas}
\makeatletter
\newlabel{more-algebra-lemma-application-formal-glueing}{{0ALK}{source}{}{sourcepointer-1}{}}
\newlabel{more-algebra-lemma-application-formal-glueing@cref}{{[originallemma][0][]0ALK}{[1][0][]source}{}{}{}}
\makeatother
\crefname{originalremark}{original remark}{original remarks}
\Crefname{originalremark}{Original remark}{Original remarks}
\makeatletter
\newlabel{more-algebra-remark-formal-glueing-algebras}{{05EU}{source}{}{sourcepointer-2}{}}
\newlabel{more-algebra-remark-formal-glueing-algebras@cref}{{[originalremark][0][]05EU}{[1][0][]source}{}{}{}}
\makeatother
\crefname{originallemma}{original lemma}{original lemmas}
\Crefname{originallemma}{Original lemma}{Original lemmas}
\makeatletter
\newlabel{algebra-lemma-nonregular-dimension-one}{{00P9}{source}{}{sourcepointer-3}{}}
\newlabel{algebra-lemma-nonregular-dimension-one@cref}{{[originallemma][0][]00P9}{[1][0][]source}{}{}{}}
\makeatother
\crefname{originallemma}{original lemma}{original lemmas}
\Crefname{originallemma}{Original lemma}{Original lemmas}
\makeatletter
\newlabel{divisors-lemma-blowing-up-gives-effective-Cartier-divisor}{{02OS}{source}{}{sourcepointer-4}{}}
\newlabel{divisors-lemma-blowing-up-gives-effective-Cartier-divisor@cref}{{[originallemma][0][]02OS}{[1][0][]source}{}{}{}}
\makeatother
\crefname{originallemma}{original lemma}{original lemmas}
\Crefname{originallemma}{Original lemma}{Original lemmas}
\makeatletter
\newlabel{divisors-lemma-blowing-up-projective}{{02NS}{source}{}{sourcepointer-5}{}}
\newlabel{divisors-lemma-blowing-up-projective@cref}{{[originallemma][0][]02NS}{[1][0][]source}{}{}{}}
\makeatother
\crefname{originallemma}{original lemma}{original lemmas}
\Crefname{originallemma}{Original lemma}{Original lemmas}
\makeatletter
\newlabel{divisors-lemma-blow-up-and-irreducible-components}{{0BFM}{source}{}{sourcepointer-6}{}}
\newlabel{divisors-lemma-blow-up-and-irreducible-components@cref}{{[originallemma][0][]0BFM}{[1][0][]source}{}{}{}}
\makeatother
\crefname{originallemma}{original lemma}{original lemmas}
\Crefname{originallemma}{Original lemma}{Original lemmas}
\makeatletter
\newlabel{varieties-lemma-finite-in-codim-1}{{0AB7}{source}{}{sourcepointer-7}{}}
\newlabel{varieties-lemma-finite-in-codim-1@cref}{{[originallemma][0][]0AB7}{[1][0][]source}{}{}{}}
\makeatother
\crefname{originallemma}{original lemma}{original lemmas}
\Crefname{originallemma}{Original lemma}{Original lemmas}
\makeatletter
\newlabel{algebra-lemma-locally-free-semi-local-free}{{02M9}{source}{}{sourcepointer-8}{}}
\newlabel{algebra-lemma-locally-free-semi-local-free@cref}{{[originallemma][0][]02M9}{[1][0][]source}{}{}{}}
\makeatother
\crefname{originallemma}{original lemma}{original lemmas}
\Crefname{originallemma}{Original lemma}{Original lemmas}
\makeatletter
\newlabel{algebra-lemma-dominate-by-dimension-1}{{00P8}{source}{}{sourcepointer-9}{}}
\newlabel{algebra-lemma-dominate-by-dimension-1@cref}{{[originallemma][0][]00P8}{[1][0][]source}{}{}{}}
\makeatother
\crefname{originallemma}{original lemma}{original lemmas}
\Crefname{originallemma}{Original lemma}{Original lemmas}
\makeatletter
\newlabel{algebra-lemma-integral-overring-surjective}{{00GQ}{source}{}{sourcepointer-10}{}}
\newlabel{algebra-lemma-integral-overring-surjective@cref}{{[originallemma][0][]00GQ}{[1][0][]source}{}{}{}}
\makeatother
\crefname{originallemma}{original lemma}{original lemmas}
\Crefname{originallemma}{Original lemma}{Original lemmas}
\makeatletter
\newlabel{algebra-lemma-surjective-spec-radical-ideal}{{00FI}{source}{}{sourcepointer-11}{}}
\newlabel{algebra-lemma-surjective-spec-radical-ideal@cref}{{[originallemma][0][]00FI}{[1][0][]source}{}{}{}}
\makeatother
\crefname{originallemma}{original lemma}{original lemmas}
\Crefname{originallemma}{Original lemma}{Original lemmas}
\makeatletter
\newlabel{algebra-lemma-pushdown-module}{{02M0}{source}{}{sourcepointer-12}{}}
\newlabel{algebra-lemma-pushdown-module@cref}{{[originallemma][0][]02M0}{[1][0][]source}{}{}{}}
\makeatother
\title{Simultaneous local factorization of nonzerodivisors after a finite extension of a 1-dimensional Noetherian local ring}
\author{The Stacks Project Authors}
\date{stacks-project:chow.tex:local-snapshot-20261011; GNU FDL 1.2 or later}
\begin{document}
\maketitle
\noindent\textit{Editorial note.} Counted proposition: Tag 0EAG, chow.tex lines 641-655 with its own proof at lines 657-752 (statement immediately followed by its proof; no non-adjacent pairing).\par\medskip
\noindent\textit{Editorial note.} The four context blocks reproduce the author's own definitions and the three supporting lemmas of the same section (Tags 02PH, 0EA6, 0EA9, 0EAE, 02Q7, 0EAF). They are supporting results, not additional counted problems.\par\medskip
\noindent\textit{Editorial note.} Every external reference is a navigation pointer to the stable Stacks tag; none of them replaces the counted construction or the length induction retained in this item.\par\medskip
\noindent\textit{Editorial note.} The retained preamble is the project's native preamble (Xy-pic setup and the theorem/macro block), not a new macro layer; documentclass amsart is supplied mechanically because the native preamble selects the project class file, which is not redistributed here.\par\medskip
\noindent\textit{Allowed standard prerequisites (editorial).}\begin{itemize}
\item Formal glueing of finite modules and algebras over a ring with finitely many maximal ideals (More on Algebra, Lemma 0ALK; Remark 05EU). Stacks tags 0ALK, 05EU. Produces the finite extension B contained in C whose localizations realize prescribed local extensions; applied by the author at chow.tex lines 681-693 and quoted again at 549-558.
\item Finiteness of the blow-up of a 1-dimensional local ring in a two-generator ideal (Divisors, Lemmas 02OS, 02NS, 0BFM; Varieties, Lemma 0AB7). Stacks tags 02OS, 02NS, 0BFM, 0AB7. Supplies R contained in R' with a = t a', b = t b' and R' = a'R' + b'R' in the second retained lemma (chow.tex lines 575-601).
\item Domination of a 1-dimensional local domain by a discrete valuation ring, and pushdown of lengths along a finite local extension. Stacks tags 00P8, 02M0 (Algebra, Lemmas 00P8 and 02M0). Bounds the exponent m in the third retained lemma (chow.tex lines 612-638) and computes length\_A(A'/xA') at chow.tex lines 716-722.
\item Comparison of multiplicities of 2-periodic complexes under a finite-length kernel and cokernel. Stacks tags 02PH, 0EA6, 0EA9 (retained context, source lines 165-180, 197-206, 282-288). Gives length\_A(A/xA) = length\_A(A'/xA') in the author's proof at chow.tex lines 713-715.
\end{itemize}\medskip
\ifdefined\hypertarget\hypertarget{sourcepointer-1}{}\fi
\noindent\textit{Source reference.} Original lemma 0ALK: sources/stacks-chow/tags.txt, source lines 13373--13373. Navigation pointer to the formal glueing lemma cited at chow.tex line 549; the counted content is the induction of lines 657-752 that applies it, not the glueing result itself.\par\medskip
\ifdefined\hypertarget\hypertarget{sourcepointer-2}{}\fi
\noindent\textit{Source reference.} Original remark 05EU: sources/stacks-chow/tags.txt, source lines 7003--7003. Remark cited at chow.tex line 553 for endowing the glued module with an algebra structure; the author's argument is unchanged and the pointer carries no proof obligation of this task.\par\medskip
\ifdefined\hypertarget\hypertarget{sourcepointer-3}{}\fi
\noindent\textit{Source reference.} Original lemma 00P9: sources/stacks-chow/tags.txt, source lines 923--923. Quoted at chow.tex lines 562-563 only to say that the following supporting lemma is very similar to it; no step of the counted proof depends on its statement.\par\medskip
\ifdefined\hypertarget\hypertarget{sourcepointer-4}{}\fi
\noindent\textit{Source reference.} Original lemma 02OS: sources/stacks-chow/tags.txt, source lines 3527--3527. Standard blow-up input cited at chow.tex lines 581-582 and 596 inside the retained supporting lemma; the pointer replaces neither that lemma nor the induction of the counted proof.\par\medskip
\ifdefined\hypertarget\hypertarget{sourcepointer-5}{}\fi
\noindent\textit{Source reference.} Original lemma 02NS: sources/stacks-chow/tags.txt, source lines 3491--3491. Projectivity of the blow-up, cited at chow.tex line 587 as a standard scheme-theoretic fact inside the retained supporting lemma.\par\medskip
\ifdefined\hypertarget\hypertarget{sourcepointer-6}{}\fi
\noindent\textit{Source reference.} Original lemma 0BFM: sources/stacks-chow/tags.txt, source lines 14390--14390. Used at chow.tex line 589 to see that every generic point of the blow-up lies over the punctured spectrum; a standard pointer inside retained context.\par\medskip
\ifdefined\hypertarget\hypertarget{sourcepointer-7}{}\fi
\noindent\textit{Source reference.} Original lemma 0AB7: sources/stacks-chow/tags.txt, source lines 13010--13010. Finiteness criterion cited at chow.tex line 590 to conclude that the blow-up is affine; a standard pointer inside retained context.\par\medskip
\ifdefined\hypertarget\hypertarget{sourcepointer-8}{}\fi
\noindent\textit{Source reference.} Original lemma 02M9: sources/stacks-chow/tags.txt, source lines 3436--3436. Cited at chow.tex line 599 to turn the locally principal ideal IR' into a principal one over a semi-local ring; standard commutative-algebra pointer.\par\medskip
\ifdefined\hypertarget\hypertarget{sourcepointer-9}{}\fi
\noindent\textit{Source reference.} Original lemma 00P8: sources/stacks-chow/tags.txt, source lines 922--922. Existence of a dominating discrete valuation ring, cited twice at chow.tex lines 615-617 and 633-634 inside the retained boundedness lemma; the pointer does not replace the counted induction.\par\medskip
\ifdefined\hypertarget\hypertarget{sourcepointer-10}{}\fi
\noindent\textit{Source reference.} Original lemma 00GQ: sources/stacks-chow/tags.txt, source lines 616--616. Surjectivity of Spec on integral extensions, cited at chow.tex line 626 inside retained context.\par\medskip
\ifdefined\hypertarget\hypertarget{sourcepointer-11}{}\fi
\noindent\textit{Source reference.} Original lemma 00FI: sources/stacks-chow/tags.txt, source lines 572--572. Surjectivity of Spec after the radical base change A to A', cited at chow.tex line 628 inside retained context.\par\medskip
\ifdefined\hypertarget\hypertarget{sourcepointer-12}{}\fi
\noindent\textit{Source reference.} Original lemma 02M0: sources/stacks-chow/tags.txt, source lines 3427--3427. Length pushdown along a finite extension, cited at chow.tex line 722 in the author's proof; it is an allowed foundational input and the length induction that uses it is retained in full.\par\medskip
\begin{definition}
\label{definition-periodic-length}
Let $(M, N, \varphi, \psi)$ be a $2$-periodic complex
over a ring $R$ whose cohomology modules have finite length.
In this case we define the {\it multiplicity} of $(M, N, \varphi, \psi)$
to be the integer
$$
e_R(M, N, \varphi, \psi) =
\text{length}_R(H^0(M, N, \varphi, \psi))
-
\text{length}_R(H^1(M, N, \varphi, \psi))
$$
In the case of a $(2, 1)$-periodic complex $(M, \varphi, \psi)$,
we denote this by $e_R(M, \varphi, \psi)$ and we will sometimes call this
the {\it (additive) Herbrand quotient}.
\end{definition}

\medskip\noindent
An example of a $(2, 1)$-periodic complex over a ring $R$ is any triple of
the form $(M, 0, \psi)$ where $M$ is an $R$-module and $\psi$ is an
$R$-linear map. If the kernel and cokernel of $\psi$ have finite length,
then we obtain
\begin{equation}
\label{equation-multiplicity-coker-ker}
e_R(M, 0, \psi) = \text{length}_R(\Coker(\psi)) - \text{length}_R(\Ker(\psi))
\end{equation}
We state and prove the obligatory lemmas on these notations.

\begin{lemma}
\label{lemma-compare-periodic-lengths}
Let $R$ be a ring. Let $f : (M, \varphi, \psi) \to (M', \varphi', \psi')$
be a map of $(2, 1)$-periodic complexes whose cohomology modules
have finite length. If $\Ker(f)$ and $\Coker(f)$ have finite length,
then $e_R(M, \varphi, \psi) = e_R(M', \varphi', \psi')$.
\end{lemma}

\section{Preparation for tame symbols}
\label{section-preparation-tame-symbol}

\noindent
In this section we put some lemma that will help us define the
tame symbol in the next section.

\begin{lemma}
\label{lemma-glue-at-max}
Let $A$ be a Noetherian ring. Let $\mathfrak m_1, \ldots, \mathfrak m_r$
be pairwise distinct maximal ideals of $A$. For $i = 1, \ldots, r$ let
$\varphi_i : A_{\mathfrak m_i} \to B_i$ be a ring map whose
kernel and cokernel are annihilated by a power
of $\mathfrak m_i$. Then there exists a ring map $\varphi : A \to B$ such
that
\begin{enumerate}
\item the localization of $\varphi$ at $\mathfrak m_i$ is
isomorphic to $\varphi_i$, and
\item $\Ker(\varphi)$ and $\Coker(\varphi)$ are annihilated
by a power of $\mathfrak m_1 \cap \ldots \cap \mathfrak m_r$.
\end{enumerate}
Moreover, if each $\varphi_i$ is finite, injective, or
surjective then so is $\varphi$.
\end{lemma}

\begin{proof}
Set $I = \mathfrak m_1 \cap \ldots \cap \mathfrak m_r$. Set
$A_i = A_{\mathfrak m_i}$ and $A' = \prod A_i$.
Then $IA' = \prod \mathfrak m_i A_i$ and $A \to A'$
is a flat ring map such that $A/I \cong A'/IA'$.
Thus we may use More on Algebra, Lemma
\ref{more-algebra-lemma-application-formal-glueing}
to see that there exists an $A$-module map $\varphi : A \to B$
with $\varphi_i$ isomorphic to the localization of $\varphi$
at $\mathfrak m_i$. Then we can use the discussion in
More on Algebra, Remark \ref{more-algebra-remark-formal-glueing-algebras}
to endow $B$ with an $A$-algebra structure
matching the given $A$-algebra structure on $B_i$.
The final statement of the lemma follows easily from
the fact that $\Ker(\varphi)_{\mathfrak m_i} \cong \Ker(\varphi_i)$
and $\Coker(\varphi)_{\mathfrak m_i} \cong \Coker(\varphi_i)$.
\end{proof}

\noindent
The following lemma is very similar to
Algebra, Lemma \ref{algebra-lemma-nonregular-dimension-one}.

\begin{lemma}
\label{lemma-Noetherian-domain-dim-1-two-elements}
Let $(R, \mathfrak m)$ be a Noetherian local ring of dimension $1$.
Let $a, b \in R$ be nonzerodivisors.
There exists a finite ring extension $R \subset R'$
with $R'/R$ annihilated by a power of $\mathfrak m$
and nonzerodivisors $t, a', b' \in R'$ such that
$a = ta'$ and $b = tb'$ and $R' = a'R' + b'R'$.
\end{lemma}

\begin{proof}
If $a$ or $b$ is a unit, then the lemma is true with $R = R'$.
Thus we may assume $a, b \in \mathfrak m$.
Set $I = (a, b)$. The idea is to blow up $R$ in $I$.
Instead of doing the algebraic argument we work geometrically.
Let $X = \text{Proj}(\bigoplus_{d \geq 0} I^d)$.
By Divisors, Lemma
\ref{divisors-lemma-blowing-up-gives-effective-Cartier-divisor}
the morphism $X \to \Spec(R)$ is an isomorphism over
the punctured spectrum $U = \Spec(R) \setminus \{\mathfrak m\}$.
Thus we may and do view $U$ as an open subscheme of $X$.
The morphism $X \to \Spec(R)$ is projective by
Divisors, Lemma \ref{divisors-lemma-blowing-up-projective}.
Also, every generic point of $X$ lies in $U$, for example
by Divisors, Lemma \ref{divisors-lemma-blow-up-and-irreducible-components}.
It follows from Varieties, Lemma \ref{varieties-lemma-finite-in-codim-1}
that $X \to \Spec(R)$ is finite. Thus $X = \Spec(R')$ is
affine and $R \to R'$ is finite. We have $R_a \cong R'_a$ as $U = D(a)$.
Hence a power of $a$ annihilates the finite $R$-module $R'/R$.
As $\mathfrak m = \sqrt{(a)}$ we see that $R'/R$ is annihilated
by a power of $\mathfrak m$. By
Divisors, Lemma \ref{divisors-lemma-blowing-up-gives-effective-Cartier-divisor}
we see that $IR'$ is a locally principal ideal.
Since $R'$ is semi-local we see that $IR'$ is principal,
see Algebra, Lemma \ref{algebra-lemma-locally-free-semi-local-free},
say $IR' = (t)$. Then we have $a = a't$ and $b = b't$ and everything is
clear.
\end{proof}

\begin{lemma}
\label{lemma-not-infinitely-divisible}
Let $(R, \mathfrak m)$ be a Noetherian local ring of dimension $1$.
Let $a, b \in R$ be nonzerodivisors with $a \in \mathfrak m$.
There exists an integer $n = n(R, a, b)$ such that for a finite ring
extension $R \subset R'$ if $b = a^m c$ for some $c \in R'$, then $m \leq n$.
\end{lemma}

\begin{proof}
Choose a minimal prime $\mathfrak q \subset R$. Observe that
$\dim(R/\mathfrak q) = 1$, in particular $R/\mathfrak q$ is not a field.
We can choose a discrete valuation ring $A$ dominating $R/\mathfrak q$
with the same fraction field, see
Algebra, Lemma \ref{algebra-lemma-dominate-by-dimension-1}. Observe that
$a$ and $b$ map to nonzero elements of $A$ as nonzerodivisors in $R$
are not contained in $\mathfrak q$. Let $v$ be the discrete valuation on $A$.
Then $v(a) > 0$ as $a \in \mathfrak m$.
We claim $n = v(b)/v(a)$ works.

\medskip\noindent
Let $R \subset R'$ be given. Set $A' = A \otimes_R R'$.
Since $\Spec(R') \to \Spec(R)$ is surjective
(Algebra, Lemma \ref{algebra-lemma-integral-overring-surjective})
also $\Spec(A') \to \Spec(A)$ is surjective
(Algebra, Lemma \ref{algebra-lemma-surjective-spec-radical-ideal}).
Pick a prime $\mathfrak q' \subset A'$ lying over $(0) \subset A$.
Then $A \subset A'' = A'/\mathfrak q'$ is a finite extension of rings
(again inducing a surjection on spectra).
Pick a maximal ideal $\mathfrak m'' \subset A''$
lying over the maximal ideal of $A$ and a discrete valuation ring
$A'''$ dominating $A''_{\mathfrak m''}$ (see lemma cited above).
Then $A \to A'''$ is an extension of discrete valuation rings
and we have $b = a^m c$ in $A'''$. Thus $v'''(b) \geq mv'''(a)$.
Since $v''' = ev$ where $e$ is the ramification index
of $A'''/A$, we find that $m \leq n$ as desired.
\end{proof}

\begin{lemma}
\label{lemma-prepare-tame-symbol}
Let $(A, \mathfrak m)$ be a Noetherian local ring of dimension $1$.
Let $r \geq 2$ and let $a_1, \ldots, a_r \in A$ be nonzerodivisors
not all units.
Then there exist
\begin{enumerate}
\item a finite ring extension $A \subset B$ with
$B/A$ annihilated by a power of $\mathfrak m$,
\item for each maximal ideal $\mathfrak m_j \subset B$
a nonzerodivisor $\pi_j \in B_j = B_{\mathfrak m_j}$, and
\item factorizations $a_i = u_{i, j} \pi_j^{e_{i, j}}$ in $B_j$
with $u_{i, j} \in B_j$ units and $e_{i, j} \geq 0$.
\end{enumerate}
\end{lemma}

\begin{proof}
Since at least one $a_i$ is not a unit we find that $\mathfrak m$
is not an associated prime of $A$. Moreover, for any $A \subset B$
as in the statement $\mathfrak m$ is not an associated prime of $B$
and $\mathfrak m_j$ is not an associate prime of $B_j$.
Keeping this in mind will help check the arguments below.

\medskip\noindent
First, we claim that it suffices to prove the lemma for $r = 2$.
We will argue this by induction on $r$; we suggest the reader
skip the proof. Suppose we are given $A \subset B$ and $\pi_j$ in
$B_j = B_{\mathfrak m_j}$ and factorizations
$a_i = u_{i, j} \pi_j^{e_{i, j}}$ for $i = 1, \ldots, r - 1$ in $B_j$
with $u_{i, j} \in B_j$ units and $e_{i, j} \geq 0$.
Then by the case $r = 2$ for $\pi_j$ and $a_r$ in $B_j$
we can find extensions $B_j \subset C_j$ and for every maximal ideal
$\mathfrak m_{j, k}$ of $C_j$ a nonzerodivisor
$\pi_{j, k} \in C_{j, k} = (C_j)_{\mathfrak m_{j, k}}$
and factorizations
$$
\pi_j = v_{j, k} \pi_{j, k}^{f_{j, k}}
\quad\text{and}\quad
a_r = w_{j, k} \pi_{j, k}^{g_{j, k}}
$$
as in the lemma. There exists a unique finite extension $B \subset C$
with $C/B$ annihilated by a power of $\mathfrak m$ such
that $C_j \cong C_{\mathfrak m_j}$ for all $j$, see
Lemma \ref{lemma-glue-at-max}.
The maximal ideals of $C$ correspond $1$-to-$1$
to the maximal ideals $\mathfrak m_{j, k}$ in the localizations
and in these localizations we have
$$
a_i = u_{i, j} \pi_j^{e_{i, j}} =
u_{i, j} v_{j, k}^{e_{i, j}} \pi_{j, k}^{e_{i, j}f_{j, k}}
$$
for $i \leq r - 1$. Since $a_r$ factors correctly too the
proof of the induction step is complete.

\medskip\noindent
Proof of the case $r = 2$. We will use induction on
$$
\ell = \min(\text{length}_A(A/a_1A),\ \text{length}_A(A/a_2A)).
$$
If $\ell = 0$, then either $a_1$ or $a_2$ is a unit and
the lemma holds with $A = B$. Thus we may and do assume $\ell > 0$.

\medskip\noindent
Suppose we have a finite extension of rings $A \subset A'$ such that
$A'/A$ is annihilated by a power of $\mathfrak m$ and such that
$\mathfrak m$ is not an associated prime of $A'$.
Let $\mathfrak m_1, \ldots, \mathfrak m_r \subset A'$
be the maximal ideals and set $A'_i = A'_{\mathfrak m_i}$.
If we can solve the problem for $a_1, a_2$ in each $A'_i$,
then we can apply Lemma \ref{lemma-glue-at-max}
to produce a solution for $a_1, a_2$ in $A$.
Choose $x \in \{a_1, a_2\}$ such that $\ell = \text{length}_A(A/xA)$.
By Lemma \ref{lemma-compare-periodic-lengths} 
and (\ref{equation-multiplicity-coker-ker})
we have $\text{length}_A(A/xA) = \text{length}_A(A'/xA')$.
On the other hand, we have
$$
\text{length}_A(A'/xA') =
\sum [\kappa(\mathfrak m_i) : \kappa(\mathfrak m)]
\text{length}_{A'_i}(A'_i/xA'_i)
$$
by Algebra, Lemma \ref{algebra-lemma-pushdown-module}.
Since $x \in \mathfrak m$ we see that each term on the right hand side
is positive. We conclude that the induction hypothesis applies
to $a_1, a_2$ in each $A'_i$ if $r > 1$ or if $r = 1$ and
$[\kappa(\mathfrak m_1) : \kappa(\mathfrak m)] > 1$.
We conclude that we may assume each $A'$ as above is local with
the same residue field as $A$.

\medskip\noindent
Applying the discussion of the previous paragraph,
we may replace $A$ by the ring constructed in
Lemma \ref{lemma-Noetherian-domain-dim-1-two-elements}
for $a_1, a_2 \in A$. Then since $A$ is local we find,
after possibly switching $a_1$ and $a_2$, that $a_2 \in (a_1)$.
Write $a_2 = a_1^m c$ with $m > 0$ maximal. In fact, by
Lemma \ref{lemma-not-infinitely-divisible}
we may assume $m$ is maximal even after replacing $A$
by any finite extension $A \subset A'$ as in the previous paragraph.
If $c$ is a unit, then we are done. If not, then we replace
$A$ by the ring constructed in
Lemma \ref{lemma-Noetherian-domain-dim-1-two-elements}
for $a_1, c \in A$. Then either (1) $c = a_1 c'$ or
(2) $a_1 = c a'_1$. The first case cannot happen since
it would give $a_2 = a_1^{m + 1} c'$ contradicting the
maximality of $m$. In the second case we get
$a_1 = c a'_1$ and $a_2 = c^{m + 1} (a'_1)^m$.
Then it suffices to prove the lemma for $A$ and $c, a'_1$.
If $a'_1$ is a unit we're done and if not, then
$\text{length}_A(A/cA) < \ell$ because $cA$ is a strictly
bigger ideal than $a_1A$. Thus we win by induction hypothesis.
\end{proof}

\end{document}
